\begin{lemma}
Let \(V\) be a finite vertex set, let \(\Delta\in\mathbb N\), and let
\(0<\gamma\le \Delta\).  Put \(p=\gamma/\Delta\).  Then the Bernoulli
activation weights
\[
w(A)=p^{|A|}(1-p)^{|V|-|A|}\qquad(A\subseteq V)
\]
are nonnegative and satisfy \(\sum_{A\subseteq V} w(A)=1\).
\end{lemma}
\begin{proof}
Since \(0<\gamma\le\Delta\), necessarily \(\Delta>0\), so \(p=\gamma/\Delta\)
is defined as a real number with \(0<p\le 1\).  Hence \(p\ge 0\) and
\(1-p\ge 0\).  For every activation set \(A\), both factors
\(p^{|A|}\) and \((1-p)^{|V|-|A|}\) are therefore nonnegative, so
\(w(A)\ge0\).

For the normalization, list the vertices of \(V\) in any order.  Expanding the
finite product
\[
\prod_{v\in V}\bigl(p+(1-p)\bigr)
\]
chooses, independently for each vertex \(v\), either the \(p\)-term or the
\((1-p)\)-term.  The choices of the \(p\)-term determine a unique subset
\(A\subseteq V\), and the corresponding monomial is exactly
\(p^{|A|}(1-p)^{|V|-|A|}\).  Conversely every subset \(A\) arises from exactly
one such choice.  Therefore the product expansion is precisely
\[
\sum_{A\subseteq V} p^{|A|}(1-p)^{|V|-|A|}.
\]
But each factor \(p+(1-p)\) equals \(1\), so the product equals \(1\).  This
proves the displayed normalization.
\end{proof}
